\documentclass[11pt]{amsart}

\usepackage[margin=1in]{geometry}
\usepackage{amsmath,amssymb,amsthm,mathtools}
\usepackage[shortlabels]{enumitem}
\usepackage{hyperref}
\usepackage{setspace}

\newtheorem{lemma}{Lemma}
\theoremstyle{remark}
\newtheorem{remark}[lemma]{Remark}

\newcommand{\Prob}{\mathbb P}
\newcommand{\R}{\mathbb R}
\newcommand{\eps}{\varepsilon}
\newcommand{\norm}[1]{\lVert #1 \rVert}

\begin{document}
\setstretch{1.2}

\title{A Diffuse-Limit Lemma}
\author{}
\date{}
\maketitle

\begin{lemma}[Diffuse limits are infinitely divisible]\label{lem:diffuse-id-limit}
Let $d\ge1$, and let $(Z_{n,i})_{n\ge1,\ i\ge1}$ be a triangular array of
$\R^d$-valued random variables.  Assume that, for each fixed $n$, the variables
$(Z_{n,i})_{i\ge1}$ are independent and that the row sum
\[
S_n:=\sum_{i\ge1} Z_{n,i}
\]
is well defined as a limit in probability of its finite partial sums.  Assume
also that the array is infinitesimal, in the sense that for every
$\eps>0$,
\[
\sup_{i\ge1}\Prob\{\norm{Z_{n,i}}>\eps\}\longrightarrow0.
\]
If $S_n$ converges in distribution to an $\R^d$-valued random variable $S$,
then the law of $S$ is infinitely divisible.
\end{lemma}

\begin{proof}
The finite-row version is classical: if
$U_{n,1},\dots,U_{n,k_n}$ are independent $\R^d$-valued random variables,
\[
\max_{1\le j\le k_n}\Prob\{\norm{U_{n,j}}>\eps\}\longrightarrow0
\qquad(\eps>0),
\]
and
\[
\sum_{j=1}^{k_n}U_{n,j}\Rightarrow U,
\]
then the law of $U$ is infinitely divisible.  This is the classical
Kolmogorov--Khintchine theorem for null arrays; see
\cite[Theorem~9.3]{sato-levy}.  In the one-dimensional case it is also the
classical theorem of Gnedenko--Kolmogorov
\cite[Chapter~IV, Section~24]{gnedenko-kolmogorov}.

It remains only to reduce the present countable-row formulation to that
finite-row theorem.  Since $S_n$ is the limit in probability of its finite
partial sums, for each $n$ we may choose an integer $k_n$ such that
\[
\Prob\left\{
\norm{S_n-\sum_{i=1}^{k_n}Z_{n,i}}>\frac1n
\right\}\le \frac1n.
\]
Set
\[
S_n^{(k)}:=\sum_{i=1}^{k_n}Z_{n,i}.
\]
Then $S_n^{(k)}-S_n\to0$ in probability.  Hence, by Slutsky's theorem,
\[
S_n^{(k)}\Rightarrow S.
\]
For each fixed $n$, the summands $Z_{n,1},\dots,Z_{n,k_n}$ are independent,
and the finite array remains infinitesimal because
\[
\max_{1\le i\le k_n}\Prob\{\norm{Z_{n,i}}>\eps\}
\le
\sup_{i\ge1}\Prob\{\norm{Z_{n,i}}>\eps\}
\longrightarrow0.
\]
Applying the finite-row theorem gives that the law of $S$ is infinitely
divisible.
\end{proof}

\begin{remark}[Use in the limit-law lemma]
In the compactness proof one applies Lemma~\ref{lem:diffuse-id-limit} with
$d=2$, $Z_{n,1}:=0$, and, for $i\ge2$,
\[
Z_{n,i}
:=
\bigl(a_i^{(n)}\xi_{n,i},\,b_i^{(n)}\xi_{n,i}\bigr)
\mathbf 1_{\{i>m_n\}}.
\]
The row sums are
\[
\sum_{i\ge1}Z_{n,i}
=
\bigl(X_n^{\mathrm{tail}},Y_n^{\mathrm{tail}}\bigr).
\]
They are well defined in $L^2$, hence in probability, because the coefficients
belong to $\ell^2$ and the row variables are centered and $L^2$-normalized.
The array is infinitesimal: for every $\eps>0$,
\[
\sup_{i\ge1}
\Prob\{\norm{Z_{n,i}}>\eps\}
\le
\frac1{\eps^2}
\sup_{i>m_n}\bigl((a_i^{(n)})^2+(b_i^{(n)})^2\bigr)
\longrightarrow0.
\]
Therefore every weak subsequential limit of the tail pairs
$\bigl(X_n^{\mathrm{tail}},Y_n^{\mathrm{tail}}\bigr)$ has an infinitely
divisible law on $\R^2$.
\end{remark}

% Suggested bibliography entry if this lemma is imported into RealSPR.tex:
%
% @book{sato-levy,
%   author={Sato, Ken-iti},
%   title={L\'evy Processes and Infinitely Divisible Distributions},
%   series={Cambridge Studies in Advanced Mathematics},
%   volume={68},
%   publisher={Cambridge University Press},
%   year={1999},
% }

\begin{thebibliography}{9}

\bibitem{gnedenko-kolmogorov}
B.~V. Gnedenko and A.~N. Kolmogorov,
\emph{Limit Distributions for Sums of Independent Random Variables},
Addison-Wesley, 1954. Translated by K.~L. Chung.

\bibitem{sato-levy}
Ken-iti Sato,
\emph{L\'evy Processes and Infinitely Divisible Distributions},
Cambridge Studies in Advanced Mathematics, vol.~68,
Cambridge University Press, 1999.

\end{thebibliography}

\end{document}
